% --- Archon physics formalization source begin ---
% archon:physics
% archon:covers PhyXMiniProblems/problem_phyx_mini_0741.lean
% archon:source-report reports/phyx_mini/problem_phyx_mini_0741.source.json

\chapter{Physics problem phyx\_mini\_0741}
\label{ch:PhyXMiniProblems_problem_phyx_mini_0741}

\paragraph{Problem source.}
\#\# Physical scenario

In figure, a ball is thrown leftward from the left edge of the roof, at height \$h\$ above the ground. The ball hits the ground \$1.50\textbackslash{} s\$ later, at distance \$d = 25.0\textbackslash{} m\$ from the building and at angle \$\textbackslash{}theta = 60.0\textasciicircum{}\textbackslash{}circ\$ with the horizontal.

\#\# Question

What is the angle relative to the horizontal of the velocity at which the ball is thrown?

\#\# Answer choices

- A: A: \$38.2\textasciicircum{}\textbackslash{}circ\$

- B: B: \$39.6\textasciicircum{}\textbackslash{}circ\$

- C: C: \$40.4\textasciicircum{}\textbackslash{}circ\$

- D: D: \$42.4\textasciicircum{}\textbackslash{}circ\$

\#\# Figure caption (auxiliary; use the image as primary evidence)

The image depicts a diagram showing a building and an inclined plane. Here's a detailed description:

- **Building**: A rectangular structure with visible bricks and three floors. Each floor has two windows, making a total of six windows.
- **Inclined Plane**: To the left of the building, a green slope is shown with an angle labeled \textbackslash{}(\textbackslash{}theta\textbackslash{}) from the horizontal ground.
- **Ground**: A horizontal line separates the building and the inclined plane, extending under both.
- **Measurements**:
  - The horizontal distance from the base of the building to the bottom of the inclined plane is labeled \textbackslash{}(d\textbackslash{}).
  - The vertical height of the building is labeled \textbackslash{}(h\textbackslash{}).

Overall, this figure seems to illustrate a physical scenario, possibly related to geometry or physics, involving an angle of elevation and a height measurement.

\#\# Dataset metadata

Kinematics; Multi-Formula Reasoning, Spatial Relation Reasoning

\paragraph{Recorded answer/context.}
C: C: \$40.4\textasciicircum{}\textbackslash{}circ\$

\paragraph{Figure/image path.}
/root/proposal\_for\_physic/science-mango/phyx\_mini\_run/phyx\_data/test\_image/741.png

\paragraph{Formalization target.}
create a compiling Lean file with sorry bodies at `PhyXMiniProblems/problem\_phyx\_mini\_0741.lean`. The Lean declarations must preserve the physical quantities, dimensions or dimensional roles, figure labels, governing-law hypotheses, and final relation expressed by this problem.
Use Mathlib/Physlib names found through LeanExplore where available. If a physics API is missing, introduce faithful local abstractions rather than scalar placeholder aliases.

\begin{theorem}[Physics formalization target]
\label{thm:physics:phyx_mini_0741:target}
\lean{PhyXMiniProblems.ProblemPhyXMini0741.problem_phyx_mini_0741}
\uses{def:physics:phyx-mini-0741:phyxminiproblems-problemphyxmini0741-roofprojectilesetup, def:physics:phyx-mini-0741:phyxminiproblems-problemphyxmini0741-acutevelocityangleradians, def:physics:phyx-mini-0741:phyxminiproblems-problemphyxmini0741-matchesroofprojectilescenario, def:physics:phyx-mini-0741:phyxminiproblems-problemphyxmini0741-matchesroofprojectilereadouts, def:physics:phyx-mini-0741:phyxminiproblems-problemphyxmini0741-usesstandardnearearthgravity, def:physics:phyx-mini-0741:phyxminiproblems-problemphyxmini0741-matchessuppliedroofprojectilefigure, def:physics:phyx-mini-0741:phyxminiproblems-problemphyxmini0741-hasphysicalroofprojectileparameters, def:physics:phyx-mini-0741:phyxminiproblems-problemphyxmini0741-satisfiesuniformgravitykinematics, def:physics:phyx-mini-0741:phyxminiproblems-problemphyxmini0741-isuniqueroundedlaunchangleanswer}
The assigned autoformalize agent should translate this physics problem into Lean declarations in the covered file, with theorem and lemma proof bodies written as `by sorry`.
\end{theorem}
\begin{proof}
This is an autoformalization task, not a proof task. Produce statements that can later be proved by the physics prover without weakening the physical model.
\end{proof}

% --- Archon declaration topology begin ---
\section{Declaration topology}
\begin{definition}[Length Magnitude]
  \label{def:physics:phyx-mini-0741:phyxminiproblems-problemphyxmini0741-lengthmagnitude}
  \lean{PhyXMiniProblems.ProblemPhyXMini0741.LengthMagnitude}
  A nonnegative physical length, independent of the chosen unit system
\end{definition}
\begin{proof}
  This is the declared model definition of Length Magnitude.
\end{proof}
\begin{definition}[Signed Length]
  \label{def:physics:phyx-mini-0741:phyxminiproblems-problemphyxmini0741-signedlength}
  \lean{PhyXMiniProblems.ProblemPhyXMini0741.SignedLength}
  A signed physical Cartesian length component
\end{definition}
\begin{proof}
  This is the declared model definition of Signed Length.
\end{proof}
\begin{definition}[Duration]
  \label{def:physics:phyx-mini-0741:phyxminiproblems-problemphyxmini0741-duration}
  \lean{PhyXMiniProblems.ProblemPhyXMini0741.Duration}
  A nonnegative physical duration
\end{definition}
\begin{proof}
  This is the declared model definition of Duration.
\end{proof}
\begin{definition}[Velocity Component]
  \label{def:physics:phyx-mini-0741:phyxminiproblems-problemphyxmini0741-velocitycomponent}
  \lean{PhyXMiniProblems.ProblemPhyXMini0741.VelocityComponent}
  A signed physical velocity component
\end{definition}
\begin{proof}
  This is the declared model definition of Velocity Component.
\end{proof}
\begin{definition}[Acceleration Magnitude]
  \label{def:physics:phyx-mini-0741:phyxminiproblems-problemphyxmini0741-accelerationmagnitude}
  \lean{PhyXMiniProblems.ProblemPhyXMini0741.AccelerationMagnitude}
  A nonnegative physical acceleration magnitude
\end{definition}
\begin{proof}
  This is the declared model definition of Acceleration Magnitude.
\end{proof}
\begin{definition}[length In Meters]
  \label{def:physics:phyx-mini-0741:phyxminiproblems-problemphyxmini0741-lengthinmeters}
  \lean{PhyXMiniProblems.ProblemPhyXMini0741.lengthInMeters}
  \uses{def:physics:phyx-mini-0741:phyxminiproblems-problemphyxmini0741-lengthmagnitude}
  Read a nonnegative physical length in SI meters
\end{definition}
\begin{proof}
  This is the declared model definition of length In Meters.
\end{proof}
\begin{definition}[signed Length In Meters]
  \label{def:physics:phyx-mini-0741:phyxminiproblems-problemphyxmini0741-signedlengthinmeters}
  \lean{PhyXMiniProblems.ProblemPhyXMini0741.signedLengthInMeters}
  \uses{def:physics:phyx-mini-0741:phyxminiproblems-problemphyxmini0741-signedlength}
  Read a signed Cartesian length component in SI meters
\end{definition}
\begin{proof}
  This is the declared model definition of signed Length In Meters.
\end{proof}
\begin{definition}[duration In Seconds]
  \label{def:physics:phyx-mini-0741:phyxminiproblems-problemphyxmini0741-durationinseconds}
  \lean{PhyXMiniProblems.ProblemPhyXMini0741.durationInSeconds}
  \uses{def:physics:phyx-mini-0741:phyxminiproblems-problemphyxmini0741-duration}
  Read a physical duration in SI seconds
\end{definition}
\begin{proof}
  This is the declared model definition of duration In Seconds.
\end{proof}
\begin{definition}[velocity In Meters Per Second]
  \label{def:physics:phyx-mini-0741:phyxminiproblems-problemphyxmini0741-velocityinmeterspersecond}
  \lean{PhyXMiniProblems.ProblemPhyXMini0741.velocityInMetersPerSecond}
  \uses{def:physics:phyx-mini-0741:phyxminiproblems-problemphyxmini0741-velocitycomponent}
  Read a signed velocity component in SI meters per second
\end{definition}
\begin{proof}
  This is the declared model definition of velocity In Meters Per Second.
\end{proof}
\begin{definition}[acceleration In Meters Per Second Squared]
  \label{def:physics:phyx-mini-0741:phyxminiproblems-problemphyxmini0741-accelerationinmeterspersecondsquared}
  \lean{PhyXMiniProblems.ProblemPhyXMini0741.accelerationInMetersPerSecondSquared}
  \uses{def:physics:phyx-mini-0741:phyxminiproblems-problemphyxmini0741-accelerationmagnitude}
  Read an acceleration magnitude in SI meters per second squared
\end{definition}
\begin{proof}
  This is the declared model definition of acceleration In Meters Per Second Squared.
\end{proof}
\begin{definition}[degrees To Radians]
  \label{def:physics:phyx-mini-0741:phyxminiproblems-problemphyxmini0741-degreestoradians}
  \lean{PhyXMiniProblems.ProblemPhyXMini0741.degreesToRadians}
  Convert a dimensionless degree readout to a radian readout
\end{definition}
\begin{proof}
  This is the declared model definition of degrees To Radians.
\end{proof}
\begin{definition}[Planar Axis]
  \label{def:physics:phyx-mini-0741:phyxminiproblems-problemphyxmini0741-planaraxis}
  \lean{PhyXMiniProblems.ProblemPhyXMini0741.PlanarAxis}
  The horizontal and vertical axes of the planar projectile diagram
\end{definition}
\begin{proof}
  This is the declared model definition of Planar Axis.
\end{proof}
\begin{definition}[Figure Object]
  \label{def:physics:phyx-mini-0741:phyxminiproblems-problemphyxmini0741-figureobject}
  \lean{PhyXMiniProblems.ProblemPhyXMini0741.FigureObject}
  Physical objects and geometrical marks visible in image 741.png
\end{definition}
\begin{proof}
  This is the declared model definition of Figure Object.
\end{proof}
\begin{definition}[Figure Label]
  \label{def:physics:phyx-mini-0741:phyxminiproblems-problemphyxmini0741-figurelabel}
  \lean{PhyXMiniProblems.ProblemPhyXMini0741.FigureLabel}
  The three symbolic labels visible in the supplied projectile figure
\end{definition}
\begin{proof}
  This is the declared model definition of Figure Label.
\end{proof}
\begin{definition}[Planar Projectile Trajectory]
  \label{def:physics:phyx-mini-0741:phyxminiproblems-problemphyxmini0741-planarprojectiletrajectory}
  \lean{PhyXMiniProblems.ProblemPhyXMini0741.PlanarProjectileTrajectory}
  \uses{def:physics:phyx-mini-0741:phyxminiproblems-problemphyxmini0741-signedlength, def:physics:phyx-mini-0741:phyxminiproblems-problemphyxmini0741-velocitycomponent, def:physics:phyx-mini-0741:phyxminiproblems-problemphyxmini0741-planaraxis}
  An independent planar trajectory. Its argument is elapsed time measured in seconds from launch; its values remain dimensionful position and velocity components rather than untyped scalar aliases
\end{definition}
\begin{proof}
  This is the declared model definition of Planar Projectile Trajectory.
\end{proof}
\begin{definition}[Supplied Roof Projectile Figure]
  \label{def:physics:phyx-mini-0741:phyxminiproblems-problemphyxmini0741-suppliedroofprojectilefigure}
  \lean{PhyXMiniProblems.ProblemPhyXMini0741.SuppliedRoofProjectileFigure}
  \uses{def:physics:phyx-mini-0741:phyxminiproblems-problemphyxmini0741-lengthmagnitude, def:physics:phyx-mini-0741:phyxminiproblems-problemphyxmini0741-figureobject, def:physics:phyx-mini-0741:phyxminiproblems-problemphyxmini0741-figurelabel}
  Literal labels and spatial relations transcribed from the primary image. The bitmap names the unknown height and the given distance and impact angle, but contains no numerical value for the requested launch angle
\end{definition}
\begin{proof}
  This is the declared model definition of Supplied Roof Projectile Figure.
\end{proof}
\begin{definition}[Roof Projectile Setup]
  \label{def:physics:phyx-mini-0741:phyxminiproblems-problemphyxmini0741-roofprojectilesetup}
  \lean{PhyXMiniProblems.ProblemPhyXMini0741.RoofProjectileSetup}
  \uses{def:physics:phyx-mini-0741:phyxminiproblems-problemphyxmini0741-lengthmagnitude, def:physics:phyx-mini-0741:phyxminiproblems-problemphyxmini0741-duration, def:physics:phyx-mini-0741:phyxminiproblems-problemphyxmini0741-accelerationmagnitude, def:physics:phyx-mini-0741:phyxminiproblems-problemphyxmini0741-planarprojectiletrajectory, def:physics:phyx-mini-0741:phyxminiproblems-problemphyxmini0741-suppliedroofprojectilefigure}
  The independent physical quantities in the roof-projectile experiment. Neither the trajectory nor any setup field is defined from an answer choice
\end{definition}
\begin{proof}
  This is the declared model definition of Roof Projectile Setup.
\end{proof}
\begin{definition}[horizontal Position Meters]
  \label{def:physics:phyx-mini-0741:phyxminiproblems-problemphyxmini0741-horizontalpositionmeters}
  \lean{PhyXMiniProblems.ProblemPhyXMini0741.horizontalPositionMeters}
  \uses{def:physics:phyx-mini-0741:phyxminiproblems-problemphyxmini0741-signedlengthinmeters, def:physics:phyx-mini-0741:phyxminiproblems-problemphyxmini0741-roofprojectilesetup}
  Horizontal position in meters at the stated elapsed time in seconds
\end{definition}
\begin{proof}
  This is the declared model definition of horizontal Position Meters.
\end{proof}
\begin{definition}[vertical Position Meters]
  \label{def:physics:phyx-mini-0741:phyxminiproblems-problemphyxmini0741-verticalpositionmeters}
  \lean{PhyXMiniProblems.ProblemPhyXMini0741.verticalPositionMeters}
  \uses{def:physics:phyx-mini-0741:phyxminiproblems-problemphyxmini0741-signedlengthinmeters, def:physics:phyx-mini-0741:phyxminiproblems-problemphyxmini0741-roofprojectilesetup}
  Vertical position in meters at the stated elapsed time in seconds
\end{definition}
\begin{proof}
  This is the declared model definition of vertical Position Meters.
\end{proof}
\begin{definition}[horizontal Velocity Meters Per Second]
  \label{def:physics:phyx-mini-0741:phyxminiproblems-problemphyxmini0741-horizontalvelocitymeterspersecond}
  \lean{PhyXMiniProblems.ProblemPhyXMini0741.horizontalVelocityMetersPerSecond}
  \uses{def:physics:phyx-mini-0741:phyxminiproblems-problemphyxmini0741-velocityinmeterspersecond, def:physics:phyx-mini-0741:phyxminiproblems-problemphyxmini0741-roofprojectilesetup}
  Horizontal velocity in meters per second at the stated elapsed time
\end{definition}
\begin{proof}
  This is the declared model definition of horizontal Velocity Meters Per Second.
\end{proof}
\begin{definition}[vertical Velocity Meters Per Second]
  \label{def:physics:phyx-mini-0741:phyxminiproblems-problemphyxmini0741-verticalvelocitymeterspersecond}
  \lean{PhyXMiniProblems.ProblemPhyXMini0741.verticalVelocityMetersPerSecond}
  \uses{def:physics:phyx-mini-0741:phyxminiproblems-problemphyxmini0741-velocityinmeterspersecond, def:physics:phyx-mini-0741:phyxminiproblems-problemphyxmini0741-roofprojectilesetup}
  Vertical velocity in meters per second at the stated elapsed time
\end{definition}
\begin{proof}
  This is the declared model definition of vertical Velocity Meters Per Second.
\end{proof}
\begin{definition}[acute Velocity Angle Radians]
  \label{def:physics:phyx-mini-0741:phyxminiproblems-problemphyxmini0741-acutevelocityangleradians}
  \lean{PhyXMiniProblems.ProblemPhyXMini0741.acuteVelocityAngleRadians}
  \uses{def:physics:phyx-mini-0741:phyxminiproblems-problemphyxmini0741-roofprojectilesetup, def:physics:phyx-mini-0741:phyxminiproblems-problemphyxmini0741-horizontalvelocitymeterspersecond, def:physics:phyx-mini-0741:phyxminiproblems-problemphyxmini0741-verticalvelocitymeterspersecond}
  The acute angle between a nonvertical velocity and the horizontal, expressed in radians. Absolute component magnitudes make this independent of whether the depicted velocity points left or right, above or below the horizontal
\end{definition}
\begin{proof}
  This is the declared model definition of acute Velocity Angle Radians.
\end{proof}
\begin{definition}[Matches Roof Projectile Scenario]
  \label{def:physics:phyx-mini-0741:phyxminiproblems-problemphyxmini0741-matchesroofprojectilescenario}
  \lean{PhyXMiniProblems.ProblemPhyXMini0741.MatchesRoofProjectileScenario}
  \uses{def:physics:phyx-mini-0741:phyxminiproblems-problemphyxmini0741-lengthinmeters, def:physics:phyx-mini-0741:phyxminiproblems-problemphyxmini0741-durationinseconds, def:physics:phyx-mini-0741:phyxminiproblems-problemphyxmini0741-roofprojectilesetup, def:physics:phyx-mini-0741:phyxminiproblems-problemphyxmini0741-horizontalpositionmeters, def:physics:phyx-mini-0741:phyxminiproblems-problemphyxmini0741-verticalpositionmeters, def:physics:phyx-mini-0741:phyxminiproblems-problemphyxmini0741-horizontalvelocitymeterspersecond, def:physics:phyx-mini-0741:phyxminiproblems-problemphyxmini0741-verticalvelocitymeterspersecond}
  The roof edge is the coordinate origin horizontally, the ground is y = 0, and the impact point lies to the left by the positive distance d
\end{definition}
\begin{proof}
  This is the declared model definition of Matches Roof Projectile Scenario.
\end{proof}
\begin{definition}[Matches Roof Projectile Readouts]
  \label{def:physics:phyx-mini-0741:phyxminiproblems-problemphyxmini0741-matchesroofprojectilereadouts}
  \lean{PhyXMiniProblems.ProblemPhyXMini0741.MatchesRoofProjectileReadouts}
  \uses{def:physics:phyx-mini-0741:phyxminiproblems-problemphyxmini0741-lengthinmeters, def:physics:phyx-mini-0741:phyxminiproblems-problemphyxmini0741-durationinseconds, def:physics:phyx-mini-0741:phyxminiproblems-problemphyxmini0741-degreestoradians, def:physics:phyx-mini-0741:phyxminiproblems-problemphyxmini0741-roofprojectilesetup, def:physics:phyx-mini-0741:phyxminiproblems-problemphyxmini0741-acutevelocityangleradians}
  Numerical data supplied by the prose: 1.50 s, 25.0 m, and an impact angle of 60.0 degrees. No initial-angle value or answer choice is included
\end{definition}
\begin{proof}
  This is the declared model definition of Matches Roof Projectile Readouts.
\end{proof}
\begin{definition}[Uses Standard Near Earth Gravity]
  \label{def:physics:phyx-mini-0741:phyxminiproblems-problemphyxmini0741-usesstandardnearearthgravity}
  \lean{PhyXMiniProblems.ProblemPhyXMini0741.UsesStandardNearEarthGravity}
  \uses{def:physics:phyx-mini-0741:phyxminiproblems-problemphyxmini0741-accelerationinmeterspersecondsquared, def:physics:phyx-mini-0741:phyxminiproblems-problemphyxmini0741-roofprojectilesetup}
  Standard near-Earth gravitational acceleration used by the problem
\end{definition}
\begin{proof}
  This is the declared model definition of Uses Standard Near Earth Gravity.
\end{proof}
\begin{definition}[Matches Supplied Roof Projectile Figure]
  \label{def:physics:phyx-mini-0741:phyxminiproblems-problemphyxmini0741-matchessuppliedroofprojectilefigure}
  \lean{PhyXMiniProblems.ProblemPhyXMini0741.MatchesSuppliedRoofProjectileFigure}
  \uses{def:physics:phyx-mini-0741:phyxminiproblems-problemphyxmini0741-roofprojectilesetup}
  Objects, labels, and orientation relations read from image 741.png
\end{definition}
\begin{proof}
  This is the declared model definition of Matches Supplied Roof Projectile Figure.
\end{proof}
\begin{definition}[Has Physical Roof Projectile Parameters]
  \label{def:physics:phyx-mini-0741:phyxminiproblems-problemphyxmini0741-hasphysicalroofprojectileparameters}
  \lean{PhyXMiniProblems.ProblemPhyXMini0741.HasPhysicalRoofProjectileParameters}
  \uses{def:physics:phyx-mini-0741:phyxminiproblems-problemphyxmini0741-lengthinmeters, def:physics:phyx-mini-0741:phyxminiproblems-problemphyxmini0741-durationinseconds, def:physics:phyx-mini-0741:phyxminiproblems-problemphyxmini0741-accelerationinmeterspersecondsquared, def:physics:phyx-mini-0741:phyxminiproblems-problemphyxmini0741-roofprojectilesetup}
  Positivity and acute-angle conditions selecting the physical branch
\end{definition}
\begin{proof}
  This is the declared model definition of Has Physical Roof Projectile Parameters.
\end{proof}
\begin{definition}[Satisfies Uniform Gravity Kinematics]
  \label{def:physics:phyx-mini-0741:phyxminiproblems-problemphyxmini0741-satisfiesuniformgravitykinematics}
  \lean{PhyXMiniProblems.ProblemPhyXMini0741.SatisfiesUniformGravityKinematics}
  \uses{def:physics:phyx-mini-0741:phyxminiproblems-problemphyxmini0741-accelerationinmeterspersecondsquared, def:physics:phyx-mini-0741:phyxminiproblems-problemphyxmini0741-roofprojectilesetup, def:physics:phyx-mini-0741:phyxminiproblems-problemphyxmini0741-horizontalpositionmeters, def:physics:phyx-mini-0741:phyxminiproblems-problemphyxmini0741-verticalpositionmeters, def:physics:phyx-mini-0741:phyxminiproblems-problemphyxmini0741-horizontalvelocitymeterspersecond, def:physics:phyx-mini-0741:phyxminiproblems-problemphyxmini0741-verticalvelocitymeterspersecond}
  Uniform downward-gravity kinematics in SI component readouts. The four relations are the general constant-acceleration laws for every elapsed time; they contain no solved initial angle or displayed answer value
\end{definition}
\begin{proof}
  This is the declared model definition of Satisfies Uniform Gravity Kinematics.
\end{proof}
\begin{definition}[Answer Choice]
  \label{def:physics:phyx-mini-0741:phyxminiproblems-problemphyxmini0741-answerchoice}
  \lean{PhyXMiniProblems.ProblemPhyXMini0741.AnswerChoice}
  Labels of the four launch-angle choices printed in the problem
\end{definition}
\begin{proof}
  This is the declared model definition of Answer Choice.
\end{proof}
\begin{definition}[displayed Launch Angle Degrees]
  \label{def:physics:phyx-mini-0741:phyxminiproblems-problemphyxmini0741-displayedlaunchangledegrees}
  \lean{PhyXMiniProblems.ProblemPhyXMini0741.displayedLaunchAngleDegrees}
  \uses{def:physics:phyx-mini-0741:phyxminiproblems-problemphyxmini0741-answerchoice}
  Degree value printed beside each answer label
\end{definition}
\begin{proof}
  This is the declared model definition of displayed Launch Angle Degrees.
\end{proof}
\begin{definition}[Rounds To Displayed Tenth Degree]
  \label{def:physics:phyx-mini-0741:phyxminiproblems-problemphyxmini0741-roundstodisplayedtenthdegree}
  \lean{PhyXMiniProblems.ProblemPhyXMini0741.RoundsToDisplayedTenthDegree}
  \uses{def:physics:phyx-mini-0741:phyxminiproblems-problemphyxmini0741-degreestoradians, def:physics:phyx-mini-0741:phyxminiproblems-problemphyxmini0741-answerchoice, def:physics:phyx-mini-0741:phyxminiproblems-problemphyxmini0741-displayedlaunchangledegrees}
  Agreement with an angle displayed to the nearest tenth of a degree
\end{definition}
\begin{proof}
  This is the declared model definition of Rounds To Displayed Tenth Degree.
\end{proof}
\begin{definition}[Is Unique Rounded Launch Angle Answer]
  \label{def:physics:phyx-mini-0741:phyxminiproblems-problemphyxmini0741-isuniqueroundedlaunchangleanswer}
  \lean{PhyXMiniProblems.ProblemPhyXMini0741.IsUniqueRoundedLaunchAngleAnswer}
  \uses{def:physics:phyx-mini-0741:phyxminiproblems-problemphyxmini0741-answerchoice, def:physics:phyx-mini-0741:phyxminiproblems-problemphyxmini0741-roundstodisplayedtenthdegree}
  A choice uniquely matches the physical angle at the displayed precision
\end{definition}
\begin{proof}
  This is the declared model definition of Is Unique Rounded Launch Angle Answer.
\end{proof}
% --- Archon declaration topology end ---
% --- Archon physics formalization source end ---
